\SourceLeaf{67}rum primum a $2$ diversum, insuperque quando $A=4$; positive autem in omnibus casibus reliquis. Theorema, quale a cel. Wilson est prolatum, sub casu priori continetur. --- \emph{Ex. gr.} pro $A=15$ productum e numeris $1,2,4,7,8,11,13,14$ est $\equiv1\,(\operatorname{mod.}15)$. Demonstrationem brevitatis gratia non adiungimus: observamus tantum, eam simili modo perfici posse ut in art. praec., excepto quod congruentia $xx\equiv1$ plures quam duas radices habere potest, quae considerationes quasdam peculiares postulant. Posset etiam demonstratio ex consideratione indicum peti, similiter ut in art. 75, si ea quae mox de modulis non primis trademus conferantur.

\Article{79}
Revertimur ad enumerationem aliarum propositionum (art. 75).

\emph{Summa omnium terminorum periodi numeri cuiusvis est $\equiv0$}, uti in ex. art. 75, $1+5+12+8=26\equiv0\,(\operatorname{mod.}13)$.

\emph{Dem.} Numerus de cuius periodo agitur, sit $=a$, atque exponens ad quem pertinet, $=t$, eritque summa terminorum omnium periodi,
\[\equiv1+a+aa+a^3+\text{etc.}+a^{t-1}\equiv\frac{a^t-1}{a-1}\,(\operatorname{mod.}p)\]
At $a^t-1\equiv0$: quare summa haec semper erit $\equiv0$ (art. 22), nisi forte $a-1$ per $p$ sit divisibilis, sive $a\equiv1$; hunc igitur casum excipere oportet, si vel unum terminum \emph{periodum} vocare velimus.

\Article{80}
\emph{Productum ex omnibus radicibus primitivis est $\equiv1$}, excepto unico casu. $p=3$; tum enim una tantum datur radix primitiva, $2$.

\emph{Demonstr.} Si radix primitiva quaecunque pro basi assumitur, indices radicum omnium primitivarum erunt numeri ad $p-1$ primi simulque ipso minores. At horum numerorum summa, \emph{i. e.} index producti ex omnibus radicibus primitivis, est $\equiv0\,(\operatorname{mod.}p-1)$ adeoque productum $\equiv1\,(\operatorname{mod.}p)$; facile enim perspicitur. si $k$ fuerit numerus ad $p-1$ primus, etiam $p-1-k$ ad $p-1$ primum fore adeoque binos numeros ad $p-1$ primos summam constituere per $p-1$ divisibilem; ($k$ autem ipsi $p-1-k$ numquam aequalis esse potest, praeter casum, $p-1=2$, sive $p=3$, quem excepimus; manifesto enim $\frac{p-1}{2}$ in omnibus reliquis casibus ad $p-1$ non est primus).

\SourceLeaf{68}
\Article{81}
\emph{Summa omnium radicum primitivarum est aut} $\equiv0$ (quando $p-1$ per \emph{quadratum} aliquod est divisibilis), \emph{aut $\equiv\pm1\,(\mathit{mod.}p)$}, (quando $p-1$ est productum e numeris primis inaequalibus; quorum multitudo si est par signum positivum, si vero impar, negativum sumendum).

\emph{Ex.} $1^0$ pro $p=13$, habentur radices primitivae $2,6,7,11$, quarum summa $26\equiv0\,(\operatorname{mod.}13)$.

$2^0$ pro $p=11$, radices primitivae sunt $2,6,7,8$ quarum summa $23\equiv+1\,(\operatorname{mod.}11)$.

$3^0$ pro $p=31$, radices primitivae sunt $3,11,12,13,17.\,21,22,24$, quarum summa, $123\equiv-1\,(\operatorname{mod.}31)$.

\emph{Demonstr.} Supra demonstravimus (art. 55, II), si $p-1$ fuerit $=a^{\alpha}b^{\beta}c^{\gamma}$ etc. (designantibus $a,b,c$ etc. numeros primos inaequales) atque $A,B,C$ etc. numeri quicunque ad exponentes $a^{\alpha},b^{\beta},c^{\gamma}$ etc. respective pertinentes, omnia producta $ABC$ etc. exhibere radices primitivas. Facile vero etiam demonstrari potest, quamvis radicem primitivam per huiusmodi productum exhiberi posse et quidem unico tantum modo\HistoricalNote{*)}{Determinentur scilicet numeri $\mathfrak{a},\mathfrak{b},\mathfrak{c}$ etc. ita, ut sit $\mathfrak{a}\equiv1\,(\operatorname{mod.}a^{\alpha})$ et $\equiv0\,(\operatorname{mod.}b^{\beta}c^{\gamma}\text{etc.})$; $\mathfrak{b}\equiv1\,(\operatorname{mod.}b^{\beta})$ et $\equiv0\,(\operatorname{mod.}a^{\alpha}c^{\gamma}\text{etc.})$ etc. (vid. art. 32), unde fiet $\mathfrak{a}+\mathfrak{b}+\mathfrak{c}+\text{etc.}\equiv1\,(\operatorname{mod.}p-1)$, (art. 19). Iam si radix primitiva quaecunque, $r$, per productum $ABC$ etc. exhiberi debet accipiatur $A\equiv r^{\mathfrak{a}},B\equiv r^{\mathfrak{b}},C\equiv r^{\mathfrak{c}}$ etc., atque pertinebunt $A$ ad exponentem $a^{\alpha}$, $B$ ad exponentem $b^{\beta}$ etc.; productumque ex omnibus $A,B,C$ etc. erit $\equiv r\,(\operatorname{mod.}p)$; denique facile perspicitur $A,B,C$ etc. alio modo determinari non posse.}.

Unde sequitur haec producta loco ipsarum radicum primitivarum accipi posse. At quoniam in his productis omnes valores ipsius $A$ cum omnibus ipsius $B$ etc. combinari oportet, omnium horum productorum summa aequalis est producto ex summa omnium valorum ipsius $A$, in summam omnium valorum ipsius $B$, in summam omnium valorum ipsius $C$ etc. uti ex doctrina combinationum notum est. Designentur omnes valores ipsorum $A$; $B$ etc., per $A,A',A''$ etc.: $B,B',B''$ etc. etc., eritque summa omnium radicum primitivarum
\[\equiv(A+A'+\text{etc.})(B+B'+\text{etc.})\text{etc.}\]
Iam dico, si exponens $\alpha$ fuerit $=1$, summam $A+A'+A''+\text{etc.}$ fore $\equiv-1\,(\operatorname{mod.}p)$, si vero $\alpha$ fuerit $>1$, summam hanc fore $\equiv0$, similiterque de reliquis $\beta,\gamma$ etc. Simulac haec erunt demonstrata, theorematis nostri veritas
mani\SourceLeaf{69}festa erit. Quando enim $p-1$ per quadratum aliquod divisibilis est, aliquis exponentium $\alpha,\beta,\gamma$ etc. unitatem superabit, adeoque aliquis factorum, quorum producto congrua est summa omnium radicum primitivarum, erit $\equiv0$, et proin etiam productum ipsum: quando vero $p-1$ per nullum quadratum dividi potest, omnes exponentes $\alpha,\beta,\gamma$ etc. erunt $=1$, unde summa omnium radicum primitivarum congrua erit producto ex tot factoribus. quorum quisque $\equiv-1$, quot habentur numeri $a,b,c$ etc., adeoque erit $\equiv\pm1$, prout horum numerorum multitudo par vel impar. Illa autem ita probantur.

$1^0$. Quando $\alpha=1$ atque $A$ numerus ad exponentem $a$ pertinens, reliqui numeri ad hunc exponentem pertinentes erunt $A^2,A^3\GaussFourDots A^{a-1}$. At
\[1+A+A^2+A^3\GaussFourDots+A^{a-1}\]
est summa periodi completae, adeoque $\equiv0$ (art. 79), quare
\[A+A^2+A^3\GaussFourDots+A^{a-1}\equiv-1\]

$2^0$. Quando autem $\alpha>1$, atque $A$ numerus ad exponentem $a^{\alpha}$ pertinens. reliqui numeri ad hunc exponentem pertinentes habebuntur, si ex his $A^2,A^3,A^4\GaussFourDots A^{a^{\alpha}-1}$ reiiciuntur $A^a,A^{2a},A^{3a}$ etc., vid. art. 53; quare summa eorum erit
\[\equiv1+A+A^2\GaussFourDots+A^{a^{\alpha}-1}-(1+A^a+A^{2a}\GaussFourDots+A^{a^{\alpha}-a})\]
\emph{i. e.} congrua differentiae duarum periodorum, adeoque $\equiv0$. \emph{Q. E. D.}

\begin{center}\emph{De modulis qui sunt numerorum primorum potestates.}\end{center}

\Article{82}
Omnia quae hactenus exposuimus innituntur suppositioni, modulum esse numerum primum. Superest ut eum quoque casum consideremus, ubi pro modulo assumitur numerus compositus. Attamen quum hic neque proprietates tam elegantes eniteant, quam in casu priori, neque ad eas inveniendas artificiis subtilibus sit opus, sed potius omnia fere per solam principiorum praecedentium applicationem erui possint. omnes minutias hic exhaurire superfluum atque taediosum foret. Breviter itaque quae huic casui cum priori sint communia quaeque propria exponemus.

\SourceLeaf{70}
\Article{83}
Propositiones artt. 45---48 generaliter iam sunt demonstratae. At prop. art. 49 ita immutari debet:

\emph{Si $f$ designat, quot numeri dentur ad $m$ primi simul ipso $m$ minores, i. e. si $f=\phi m$ (art. 38): exponens $t$ infimae potestatis numeri dati $a$ ad $m$ primi, quae secundum modulum $m$ unitati est congrua, vel erit $=f$ vel pars aliquota huius numeri.}

Demonstratio prop. art. 49 etiam pro hoc casu valere potest, si modo ubique loco ipsius $p$, $m$, loco ipsius $p-1$, $f$, et loco numerorum $1,2,3,\GaussFourDots p-1$, numeri ad $m$ primi simulque ipso $m$ minores substituantur. Huc itaque lectorem ablegamus. Ceterum demonstrationes reliquae de quibus illic locuti sumus (artt. 50, 51) non sine multis ambagibus ad hunc casum applicari possunt. --- At respectu propositionum sequentium, art. 52 sqq. magna differentia incipit inter modulos, qui numerorum primorum sunt potestates, eosque, qui per plures numeros primos dividi possunt. Seorsim itaque modulos prioris generis contemplabimur.

\Article{84}
Si modulus $m=p^n$, designante $p$ numerum primum, erit $f=p^{n-1}(p-1)$ (art. 38). Iam si disquisitiones in artt. 53, 54 contentae ad hunc casum applicantur, mutatis mutandis uti in art. praec. praescripsimus, invenietur, omnia quae ibi demonstrata sunt etiam pro hoc casu locum habere, si modo ante probatum esset, congruentiam formae $x^t-1\equiv0\,(\operatorname{mod.}p^n)$ plures quam $t$ radices diversas habere non posse. Pro modulo primo hanc veritatem ex propositione generaliori art. 43 deduximus, quae autem in omni sua extensione de modulis primis tantummodo valet, neque adeo ad hunc casum applicanda. Attamen propositionem pro hoc casu particulari veram esse per methodum singularem demonstrabimus. Infra (sect. VIII) idem facilius invenire docebimus.

\Article{85}
Demonstrandum proponimus nobis hoc theorema:

\emph{Si numerorum $t$ et $p^{n-1}(p-1)$ divisor communis maximus est $e$, congruentia $x^t\equiv1\,(\mathit{mod.}p^n)$ habebit $e$ radices diversas.}

Sit $e=kp^{\nu}$ ita ut $k$ factorem $p$ non involvat, adeoque numerum $p-1$
